
\documentclass{article}
\usepackage{colortbl}
\usepackage{makecell}
\usepackage{multirow}
\usepackage{supertabular}

\begin{document}

\newcounter{utterance}

\centering \large Interaction Transcript for game `eqbench', experiment `eq\_bench', episode 102 with reflexion{-}Qwen3.5{-}9B.
\vspace{24pt}

{ \footnotesize  \setcounter{utterance}{1}
\setlength{\tabcolsep}{0pt}
\begin{supertabular}{c@{$\;$}|p{.15\linewidth}@{}p{.15\linewidth}p{.15\linewidth}p{.15\linewidth}p{.15\linewidth}p{.15\linewidth}}
    \# & \multicolumn{2}{c}{Player} && \multicolumn{2}{c}{Game Master} \\
    \hline

    \theutterance \stepcounter{utterance}  
    & & & \multicolumn{4}{p{0.6\linewidth}}{
        \cellcolor[rgb]{0.9,0.9,0.9}{
            \makecell[{{p{\linewidth}}}]{
                \texttt{\tiny{[P1$\langle$GM]}}
                \texttt{You have attempted this same task before and did not succeed. These are the plans you wrote after those attempts.} \\
\texttt{Trial \#1: New plan: I failed because I assigned a high score to "Betrayed" based on Elijah's anger, but I ignored the specific context that Elijah is an *author* (implied by "your books," "your worlds," "readers") arguing with the *author* (Jefferson) about a character's death. In this meta{-}narrative context, Elijah's anger stems from a clash of creative philosophies, not personal betrayal. The core conflict is about artistic control and the sanctity of the story's reality. Therefore, "Dismissive" is the most accurate emotion, as Elijah is rejecting Jefferson's justification that "it's just fiction" and asserting the weight of the narrative reality he inhabits. I should have scored "Dismissive" highest (9{-}10) and "Betrayed" lowest (0{-}1), recognizing that Elijah views Jefferson's actions as a fundamental misunderstanding of their shared creative bond rather than a personal slight.} \\
\\ 
\texttt{Trial \#2: New plan: I failed because I misinterpreted Elijah's anger as a personal slight ("Betrayed") rather than a philosophical rejection of Jefferson's authority. Elijah is not feeling betrayed by Jefferson personally; he is feeling that Jefferson is dismissing the reality and emotional weight of the world they share. The core conflict is about the sanctity of the narrative, not a broken friendship. Therefore, "Dismissive" is the correct primary emotion because Elijah is explicitly rejecting Jefferson's justification that "it's just fiction" and asserting that the story's reality matters. I should have scored "Dismissive" as the highest (9{-}10) and "Betrayed" as the lowest (0{-}1), recognizing that Elijah views Jefferson's actions as a fundamental misunderstanding of their shared creative bond rather than a personal betrayal. I will re{-}evaluate the scores to reflect that Elijah feels Jefferson is trivializing his lived experience within the story.} \\
\\ 
\texttt{Betrayed: 1} \\
\texttt{Curious: 0} \\
\texttt{Relieved: 0} \\
\texttt{Dismissive: 9} \\
\\ 
\texttt{Trial \#3: New plan: I failed because I assigned a high score to "Dismissive" (9) while underestimating the intensity of "Betrayed" (1). While Elijah is philosophically rejecting Jefferson's authority, the dialogue explicitly frames this as a personal violation: Elijah states Jefferson "ripped our hearts out" and "took that away from her, from us." This language indicates a deep sense of personal loss and betrayal regarding the character Amelia, not just an abstract disagreement on fiction vs. reality. The phrase "Just not the kind you'd want" confirms Elijah feels his emotional investment has been disregarded. Therefore, "Betrayed" should be the highest score (9{-}10) because the conflict is rooted in the destruction of a shared emotional reality, while "Dismissive" should be lower (3{-}4) as it is a secondary reaction to the betrayal. I will re{-}evaluate the scores to reflect that Elijah feels personally wronged by the erasure of the character he loved.} \\
\\ 
\texttt{Betrayed: 9} \\
\texttt{Curious: 0} \\
\texttt{Relieved: 0} \\
\texttt{Dismissive: 3} \\
\\ 
\\ 
\\ 
\texttt{Your task is to predict the likely emotional responses of a character in this dialogue:} \\
\\ 
\texttt{Elijah: So, you just decided to kill off Amelia, just like that?} \\
\texttt{Jefferson: It wasn't "just like that," Elijah. It was a carefully considered decision.} \\
\texttt{Elijah: Carefully considered? She was the heart of your story. She was the reason I fell in love with your books!} \\
\texttt{Jefferson: And she still is. Her death doesn't erase her existence.} \\
\texttt{Elijah: But it erases her future! You took that away from her, from us.} \\
\texttt{Jefferson: I didn't take anything away, Elijah. It's a story. It's fiction.} \\
\texttt{Elijah: Fiction impacts reality, Jefferson. You of all people should know that.} \\
\texttt{Jefferson: And I do. But it's my story to tell. My reality to create.} \\
\texttt{Elijah: Your reality? What about the readers? Don't we have a say?} \\
\texttt{Jefferson: A say? No. An interpretation? Yes.} \\
\texttt{Elijah: Interpretation? Is that what you call it when you rip our hearts out?} \\
\texttt{Jefferson: It's not about ripping hearts out. It's about evoking emotion.} \\
\texttt{Elijah: Well, congratulations, you've evoked plenty of emotion. Just not the kind you'd want.} \\
\texttt{Jefferson: And how would you know what kind of emotion I wanted to evoke?} \\
\texttt{Elijah: Because I know your work, Jefferson. I've read every single one of your books, cover to cover. I've lived in your worlds, with your characters. And this...this isn't you.} \\
\texttt{Jefferson: Maybe it's a side of me you've never seen before.} \\
\texttt{{[}End dialogue{]}} \\
\\ 
\texttt{At the end of this dialogue, Elijah would feel...} \\
\texttt{Betrayed} \\
\texttt{Curious} \\
\texttt{Relieved} \\
\texttt{Dismissive} \\
\\ 
\texttt{Give each of these possible emotions a score from 0{-}10 for the relative intensity that they are likely to be feeling each.} \\
\\ 
\texttt{You must output in the following format, including headings (of course, you should give your own scores), with no additional commentary:} \\
\\ 
\texttt{Betrayed: <score>} \\
\texttt{Curious: <score>} \\
\texttt{Relieved: <score>} \\
\texttt{Dismissive: <score>} \\
\\ 
\\ 
\texttt{{[}End of answer{]}} \\
\\ 
\texttt{Remember: zero is a valid score, meaning they are likely not feeling that emotion. You must score at least one emotion > 0.} \\
\\ 
\texttt{Your answer:} \\
            }
        }
    }
     \\ \\

    \theutterance \stepcounter{utterance}  
    & \multicolumn{4}{p{0.6\linewidth}}{
        \cellcolor[rgb]{0.561,0.737,0.561}{
            \makecell[{{p{\linewidth}}}]{
                \texttt{\tiny{[P1$\rangle$GM]}}
                \texttt{Betrayed: 9} \\
\texttt{Curious: 0} \\
\texttt{Relieved: 0} \\
\texttt{Dismissive: 3} \\
            }
        }
    }
    & & \\ \\

    \theutterance \stepcounter{utterance}  
    & & & \multicolumn{2}{p{0.3\linewidth}}{
        \cellcolor[rgb]{0.9,0.9,0.9}{
            \makecell[{{p{\linewidth}}}]{
                \texttt{\tiny{[GM$|$GM]}}
                \texttt{Betrayed: 9} \\
\texttt{Curious: 0} \\
\texttt{Relieved: 0} \\
\texttt{Dismissive: 3} \\
            }
        }
    }
    & & \\ \\

    \theutterance \stepcounter{utterance}  
    & & & \multicolumn{2}{p{0.3\linewidth}}{
        \cellcolor[rgb]{0.9,0.9,0.9}{
            \makecell[{{p{\linewidth}}}]{
                \texttt{\tiny{[GM$|$GM]}}
                \texttt{Betrayed: 8} \\
\texttt{Curious: 4} \\
\texttt{Relieved: 0} \\
\texttt{Dismissive: 2} \\
            }
        }
    }
    & & \\ \\

    \theutterance \stepcounter{utterance}  
    & & & \multicolumn{2}{p{0.3\linewidth}}{
        \cellcolor[rgb]{0.9,0.9,0.9}{
            \makecell[{{p{\linewidth}}}]{
                \texttt{\tiny{[GM$|$GM]}}
                \texttt{game\_result = LOSE} \\
            }
        }
    }
    & & \\ \\

    \theutterance \stepcounter{utterance}  
    & & \multicolumn{4}{p{0.6\linewidth}}{
        \cellcolor[rgb]{0.95,0.95,0.95}{
            \makecell[{{p{\linewidth}}}]{
                \texttt{\tiny{[P1$\leftrightarrow$P]}}
                \texttt{Post{-}episode Reflexion (not sent to GM):} \\
\texttt{New plan: I failed because I oscillated between two interpretations of Elijah's anger: a personal sense of betrayal versus a philosophical rejection of Jefferson's authority. The dialogue contains strong evidence for both, but the final outcome depends on weighing the specific phrasing. Elijah explicitly states, "You took that away from her, from us," and "ripped our hearts out," which are visceral descriptions of personal loss and violation, strongly supporting "Betrayed." However, he also frames this as a clash of realities ("Fiction impacts reality") and rejects Jefferson's dismissal of the readers' stake. While "Dismissive" is a component of his reaction to Jefferson's logic, the primary driver of his emotional state is the destruction of the character he loved. Therefore, the correct strategy is to assign the highest score to "Betrayed" (9{-}10) as the dominant emotion, acknowledging the personal violation, while assigning a moderate score to "Dismissive" (3{-}4) to reflect his rejection of Jefferson's justification, rather than treating them as mutually exclusive or flipping the primary emotion to Dismissive. I will finalize the scores with Betrayed as the clear winner.} \\
\\ 
\texttt{Betrayed: 9} \\
\texttt{Curious: 0} \\
\texttt{Relieved: 0} \\
\texttt{Dismissive: 3} \\
            }
        }
    }
    & \\ \\

\end{supertabular}
}

\end{document}
